% Separate continuation. Shared drivers and earlier packets are root-owned.
\subsection{Computational Expressiveness: Extraction, Coding, and Consumption}
\label{sec:p09-recombination-computation}

The expressive power of splicing depends on what is observed and on how
inputs are supplied. Terminal filtering, a letter-to-letter coding, and a
consuming multiset are three different interfaces. The following constructions
follow the published chapter's Lemmas~3.15--3.16 and Theorems~3.5--3.6
\cite[pp.~313--323]{HeadPaunPixton1997RecombinationIntake}.
Here \(\lambda\) is the source's empty word and \(|w|\) is word length.
The rule separators \(\#,\$\) and the oriented relation
\((x,y)\vdash_r z\) retain Definition~\ref{def:p09-recombination-splicing}.
\(RE\) denotes recursively enumerable languages. A terminal filter is an
intersection with \(T^*\); it is not a coding or a modification of a rule.

\paragraph{Finite seeds and a regular rule language.}

\begin{hbdefinition}{def:p09-recombination-regular-simulation}{The published grammar-simulation construction}
Let \(G=(N,T,S,P)\) be a type-0 grammar, with disjoint finite nonterminal
and terminal alphabets, start symbol \(S\in N\), and finitely many
productions \(u\to v\). Each left side is nonempty and contains a
nonterminal. Put \(U=N\cup T\cup\{B\}\). The symbols
\(X,X',B,Y,Z\), and one symbol \(Y_\alpha\) for each \(\alpha\in U\),
are fresh. In particular \(Y_\alpha\) is a single tagged symbol, not
the two-letter word \(Y\alpha\). Use
\[
 V=N\cup T\cup\{X,X',B,Y,Z\}\cup\{Y_\alpha:\alpha\in U\}.
\]
The seven source rule groups are the following words, with exactly the
indicated ranges:
\begin{enumerate}
\item \(Xw\#uY\$Z\#vY\), for \(u\to v\in P\), \(w\in U^*\).
\item \(Xw\#\alpha Y\$Z\#Y_\alpha\), for \(\alpha\in U\), \(w\in U^*\).
\item \(X'\alpha\#Z\$X\#wY_\alpha\), for \(\alpha\in U\), \(w\in U^*\).
\item \(X'w\#Y_\alpha\$Z\#Y\), for \(\alpha\in U\), \(w\in U^*\).
\item \(X\#Z\$X'\#wY\), for \(w\in U^*\).
\item \(\#ZY\$XB\#wY\), for \(w\in T^*\).
\item \(\#Y\$XZ\#\).
\end{enumerate}
Let \(R\) be their union and \(\sigma=(V,R)\). The finite seed is
\[
\begin{split}
 L_0={}&\{XBSY,ZY,XZ\}\cup\{ZvY:u\to v\in P\}\\
      &\cup\{ZY_\alpha,X'\alpha Z:\alpha\in U\}.
\end{split}
\]
The use of \(w\in U^*\) or \(T^*\) makes these regular rule languages,
not finite rule sets.
\end{hbdefinition}

\begin{hbtheorem}{thm:p09-recombination-regular-simulation}{proved=annex:p09.recombination-regular-simulation}{Exact terminal extraction from the regular-rule construction}
For the construction above,
\[
 \sigma^*(L_0)\cap T^*=L(G).
\]
Thus every recursively enumerable language is a terminal-filtered language
generated from a finite seed by an H scheme with a regular rule language.
\end{hbtheorem}

An ordinary active word is \(XpBqY\), with \(p,q\in(N\cup T)^*\).
Its decoded sentential form is \(qp\), not \(pBq\).
Groups 2--5 move the last symbol to the front through the four phases
\[
 Xw\alpha Y,\quad XwY_\alpha,\quad
 X'\alpha wY_\alpha,\quad X'\alpha wY,\quad X\alpha wY.
\]
The marker \(B\) determines the unchanged beginning of the decoded word,
even when it is the symbol moved. Group 1 rewrites a suffix \(u\) by \(v\).
Rotation makes an arbitrary occurrence of \(u\) into that suffix without
allowing a grammar production to cross \(B\).
Extraction first removes \(XB\) from \(XBwY\), \(w\in T^*\), using
group 6; group 7 then removes \(Y\).
Premature removal of \(Y\) leaves an inactive marked word, not a terminal
computation. The proof also treats inert \(Z\)-bearing by-products, rather
than assuming that all generated words lack \(Z\).

\paragraph{Why terminal filtering is not an unrestricted language equality.}

\begin{hbdefinition}{def:p09-recombination-wrapper}{The source's wrapper and restricted erasure}
For \(L\subseteq V^*\) and distinct \(c,d\notin V\), put
\[
 L_1=(dc)^*L(dc)^*,\qquad
 L_2=c(dc)^*L(dc)^*d,\qquad L'=L_1\cup L_2.
\]
Define the monoid morphism \(h\) by \(h(a)=a\) for \(a\in V\),
\(h(c)=h(d)=\lambda\). Its restriction to \(cV^*d\) erases exactly
two endpoint symbols. The closure used here is restricted-morphism closure
for this bounded erasure, not a tacit assumption of arbitrary erasing
morphism closure.
\end{hbdefinition}

\begin{hbtheorem}{thm:p09-recombination-wrapper-separation}{proved=annex:p09.recombination-wrapper-separation}{General wrapper obstruction}
Suppose a family \(FA\) is closed under intersection with regular
languages and under the restricted endpoint morphism just specified.
If \(L\notin FA\), then \(L_1,L_2,L'\notin FA\), and no H scheme,
with a rule set of any complexity, generates \(L'\) by ordinary iterated
set splicing from a seed language in \(FA\).
In particular \(L'\notin H(FA,RE)\).
\end{hbtheorem}

The reconstruction does not claim that a rule merely chooses an unfamiliar
prefix. On a word of \(L_1\), prefixing \(c\) and suffixing \(d\) gives
another available word of \(L_2\), with the same recognition site.
On a word of \(L_2\), the analogous wrapper is \(d\) on the left and
\(c\) on the right. The identical splice therefore also yields \(cz\)
or \(dz\), outside \(L'\). This contradicts closure of an exactly
generated language. The argument covers \(\lambda\in L\), empty
handles, and rules with arbitrary rather than computable membership.

\paragraph{Radius one after a letter-to-letter coding.}

\begin{hbdefinition}{def:p09-recombination-transition-coding}{Transition letters and the no-triples seed}
Let a regular language's nonempty part be generated by a right-linear
grammar \(G=(N,T,S,P)\) whose productions are \(X\to aY\) and
\(X\to a\), with \(a\in T\). A symbol \([X,a,Y]\) denotes an entire
transition letter; \([X,a,*]\) denotes an entire final-transition letter.
Neither is a word of three alphabet letters. Let \(V\) be the finite
set of these letters supplied by the productions.
Use a rule \(A\#\$\#C\) whenever
\(A=[X,a,Y]\) and \(C=[Y,b,Z]\) or \([Y,b,*]\) belong to \(V\).
Let \(L_0\) consist of all successful transition paths beginning at
\(S\) in which no nonfinal transition letter occurs three times.
The coding \(h:V\to T\) sends \([X,a,Y]\) and \([X,a,*]\) to \(a\).
Its extension to words is letter-to-letter. If literal empty-word equality
is required, also put \(\lambda\) in the seed exactly when it belongs
to the target language.
\end{hbdefinition}

\begin{hbtheorem}{thm:p09-recombination-transition-coding}{proved=annex:p09.recombination-transition-coding}{Exact coding by radius-one finite-seed splicing}
The seed in Definition~\ref{def:p09-recombination-transition-coding} is
finite. Its H scheme has radius at most one, and its iterated closure is
exactly the set of successful transition paths, with the stated optional
empty word. Consequently its coding is exactly the target regular language.
\end{hbtheorem}

The source states its class equality under the convention of ignoring
empty-word differences \cite[Theorem~3.5, pp.~317--318]{HeadPaunPixton1997RecombinationIntake}.
The optional seed word supplies literal equality without pretending that a
nonerasing coding sends a nonempty path to \(\lambda\).
This theorem is compatible with the strict radius hierarchy: the transition
alphabet and observation map have changed. It does not assert that every
regular language itself lies in \(H(FIN,1)\).

\paragraph{A consuming finite-rule computation.}

For multisets use the transition of the published section 3.3: two available
occurrences are consumed and both complementary products are released.
The source glyph \(M\Longrightarrow_\sigma M'\) denotes one such event;
\(\Longrightarrow_\sigma^*\) is its reflexive transitive closure.
Here \(M:V^*\to\mathbb N\cup\{\infty\}\), \(\operatorname{supp}(M)\)
is its support, and \(\sigma^*(M)\) consists of words appearing in any
reachable multiset. These uses of \(M\) and \(\sigma^*\) are distinct
from the reusable-set operation.

\begin{hbdefinition}{def:p09-recombination-consuming-simulation}{Published guards and the repaired finite construction}
Let \(G=(N,T,S,P)\) be as above, with \(1\leq|u|\leq2\) and
\(0\leq|v|\leq2\), omit identity productions, and label
the remaining productions injectively by a finite set \(K\).
Write \(r:u\to v\) and put \(U=N\cup T\).
The symbols \(X_1,X_2,Y,Z_1,Z_2,(r),[r]\) are fresh, with one
single-symbol pair \((r),[r]\) per label; the parentheses and brackets
are part of the printed names, not separate letters.
Thus
\[
 V=U\cup\{X_1,X_2,Y,Z_1,Z_2\}\cup\{(r),[r]:r\in K\}.
\]
Use the following finite rule groups. In groups 1 and 3,
\(\delta_1,\delta_2\in U\cup\{X_1\}\) and
\(\beta_1,\beta_2\in U\cup\{X_2\}\).
In groups 4--6, \(\delta\in U\cup\{X_1\}\) and
\(\beta\in U\cup\{X_2\}\).
\begin{enumerate}
\item \(\delta_1\delta_2Yu\#\beta_1\beta_2\$(r)v\#[r]\), for \(r:u\to v\in P\).
\item \(Y\#u[r]\$(r)\#v\alpha\), for \(r:u\to v\in P\), \(\alpha\in U\cup\{X_2\}\).
\item \(\delta_1\delta_2Y\alpha\#\beta_1\beta_2\$Z_1\alpha Y\#Z_2\), for \(\alpha\in U\).
\item \(\delta\#Y\alpha Z_2\$Z_1\#\alpha Y\beta\), for \(\alpha\in U\).
\item \(\delta\alpha Y\#\beta_1\beta_2\beta_3\$Z_1Y\alpha\#Z_2\), for \(\alpha,\beta_1\in U\), \(\beta_2,\beta_3\in U\cup\{X_2\}\), together with
 \(\delta\alpha Y\#X_2^2\$Z_1Y\alpha\#Z_2\), for \(\alpha\in U\).
\item \(\delta\#\alpha YZ_2\$Z_1\#Y\alpha\beta\), for \(\alpha\in U\).
\item \(\#YY\$X_1^2Y\#w\), for
 \(w\in\{X_2^2\}\cup T\{X_2^2\}\cup T^2\{X_2^2\}\cup T^3\).
\item \(\#X_2^2\$Y^3\#\).
\end{enumerate}
Group 5 here is the explicit repaired variant. The published group is only
\(\delta\alpha Y\#\beta_1\beta_2\beta_3\$Z_1Y\alpha\#Z_2\),
with \(\beta_1\in U\) and \(\beta_2,\beta_3\in U\cup\{X_2\}\).
The repair adds only the exact right-boundary case \(X_2^2\), not a
general two-letter guard. The other groups retain the published words and
ranges. In particular a two-letter grammar fragment before a tag is not
mistaken for this boundary.

Initially \(M(X_1^2YSX_2^2)=1\). Each reservoir word
\[
 (r)v[r]\ (r\in K),\qquad
 Z_1\alpha YZ_2,\quad Z_1Y\alpha Z_2\ (\alpha\in U),\qquad YY
\]
has infinite multiplicity, and every other word has multiplicity zero.
\end{hbdefinition}

The published statement uses \(1\leq|u|\leq2\), \(0\leq|v|\leq2\)
and concludes universality via a grammar normal form
\cite[Theorem~3.6, pp.~320--322]{HeadPaunPixton1997RecombinationIntake}.
Its group-5 guard cannot move \(Y\) left in a word ending
\(\alpha YX_2^2\): there are only two letters after \(Y\), neither
in \(U\). The repaired construction removes this obstruction.
The bounded-production hypothesis is retained. It ensures that an unfinished
grammar fragment cannot contain a new cursor-cut site: the tag is encountered
before a complete allowed guard. Arbitrary-length productions require a
separate construction or a grammar normalization, not an unchanged use of
these eight groups.

\begin{hbtheorem}{thm:p09-recombination-consuming-simulation}{proved=annex:p09.recombination-consuming-simulation}{Exact grammar simulation by a repaired consuming scheme}
For Definition~\ref{def:p09-recombination-consuming-simulation},
\[
 \sigma^*(M)\cap T^*=L(G).
\]
The seed support and rule set are finite, and the repaired scheme has radius
at most five. This theorem proves the complete simulation of the specified
bounded-production grammar; grammar normalization is a distinct dependency.
\end{hbtheorem}

\begin{hbcorollary}{thm:p09-recombination-consuming-universality}{cited=HeadPaunPixton1997RecombinationIntake}{The source's universal language-class claim}
Every recursively enumerable language is a terminal-filtered consuming
finite-support, finite-rule splicing language. To obtain this class statement
from the preceding simulation, one also uses a type-0 grammar normal form
with \(1\leq|u|\leq2\) and \(0\leq|v|\leq2\).
\end{hbcorollary}

The cited status of this class statement does not attribute a proof of that
grammar-normalization theorem to the preceding splicing annex.

There is only one computation, not indefinitely many copies of an active
sentential form. At a completed phase its active word is
\(X_1^2pYqX_2^2\), with multiplicity one, and its decoded form is
\(pq\). A rewriting event temporarily separates it into
\(X_1^2pYu[r]\) and \((r)vqX_2^2\), each with multiplicity one.
The matching reassembly consumes both, produces
\(X_1^2pYvqX_2^2\), and releases the inert waste \((r)u[r]\).
Right and left cursor moves likewise use separate cut/rejoin phases.
The annex proves the phase invariant, including harmless stuttering events
that a waste word can participate in, both cursor directions at the ends,
and the two-stage terminal cleanup. Reusable-set availability would permit
cross-combinations of fragments from different histories and does not
satisfy that invariant.

\paragraph{Finite populations and the zero-length word.}

\begin{hbdefinition}{def:p09-recombination-finite-population}{Population size, weight, and the reachable-state graph}
For a finite multiset with finite multiplicities define
\[
 \operatorname{cnt}(M)=\sum_xM(x),\qquad
 \operatorname{wht}(M)=\sum_x|x|M(x).
\]
For integers \(n,W\), let \(U_{n,W}\) be all finite multisets over
\(V^*\) with count \(n\) and weight \(W\). The graph
\(\Gamma(\sigma,M)\) has vertex set
\(U_{\operatorname{cnt}(M),\operatorname{wht}(M)}\), with an edge
\(M'\to M''\) exactly when \(M'\Longrightarrow_\sigma M''\).
The count condition remains necessary when \(\lambda\) is permitted.
\end{hbdefinition}

\begin{hbtheorem}{thm:p09-recombination-finite-population}{proved=annex:p09.recombination-finite-population}{Conservation and finite reachable-state computation}
Every consuming event preserves population size and total word length.
The graph of Definition~\ref{def:p09-recombination-finite-population} is
finite. For finite \(R\), or a supplied decision procedure for rule
membership, it can be constructed effectively. Reachability, the existence
of an infinite evolution, and maximal reachable outdegree are then decidable.
\end{hbtheorem}

Weight alone is insufficient to make the set of all equal-weight multisets
finite: arbitrarily many empty words have weight zero. Also, an arbitrary
uneffective rule set is not an algorithmic input. The conserved count and
explicit effectiveness hypothesis give the actual finite graph needed for
the source's trajectory questions \cite[p.~323]{HeadPaunPixton1997RecombinationIntake}.
Infinite reservoir multiplicities in the universal construction do not
satisfy the finite-population hypothesis.
